\documentclass[10pt,a4paper]{amsart}
\usepackage{amsmath,amssymb,amsthm,graphicx}
% For commutative diagrams we use Xy-pic
\usepackage[all]{xy}

% We use 2cell for 2-commutative diagrams.
\xyoption{2cell}
% Theorem environments.
%
\theoremstyle{plain}
\newtheorem{theorem}[subsection]{Theorem}
\newtheorem{proposition}[subsection]{Proposition}
\newtheorem{lemma}[subsection]{Lemma}

\theoremstyle{definition}
\newtheorem{definition}[subsection]{Definition}
\newtheorem{example}[subsection]{Example}
\newtheorem{exercise}[subsection]{Exercise}
\newtheorem{situation}[subsection]{Situation}

\theoremstyle{remark}
\newtheorem{remark}[subsection]{Remark}
\newtheorem{remarks}[subsection]{Remarks}

\numberwithin{equation}{subsection}

% Macros
%
\def\lim{\mathop{\mathrm{lim}}\nolimits}
\def\colim{\mathop{\mathrm{colim}}\nolimits}
\def\Spec{\mathop{\mathrm{Spec}}}
\def\Hom{\mathop{\mathrm{Hom}}\nolimits}
\def\Ext{\mathop{\mathrm{Ext}}\nolimits}
\def\SheafHom{\mathop{\mathcal{H}\!\mathit{om}}\nolimits}
\def\SheafExt{\mathop{\mathcal{E}\!\mathit{xt}}\nolimits}
\def\Sch{\mathit{Sch}}
\def\Mor{\mathop{\mathrm{Mor}}\nolimits}
\def\Ob{\mathop{\mathrm{Ob}}\nolimits}
\def\Sh{\mathop{\mathit{Sh}}\nolimits}
\def\NL{\mathop{N\!L}\nolimits}
\def\CH{\mathop{\mathrm{CH}}\nolimits}
\def\proetale{{pro\text{-}\acute{e}tale}}
\def\etale{{\acute{e}tale}}
\def\QCoh{\mathit{QCoh}}
\def\Ker{\mathop{\mathrm{Ker}}}
\def\Im{\mathop{\mathrm{Im}}}
\def\Coker{\mathop{\mathrm{Coker}}}
\def\Coim{\mathop{\mathrm{Coim}}}

% Boxtimes
%
\DeclareMathSymbol{\boxtimes}{\mathbin}{AMSa}{"02}

%
% Macros for moduli stacks/spaces
%
\def\QCohstack{\mathcal{QC}\!\mathit{oh}}
\def\Cohstack{\mathcal{C}\!\mathit{oh}}
\def\Spacesstack{\mathcal{S}\!\mathit{paces}}
\def\Quotfunctor{\mathrm{Quot}}
\def\Hilbfunctor{\mathrm{Hilb}}
\def\Curvesstack{\mathcal{C}\!\mathit{urves}}
\def\Polarizedstack{\mathcal{P}\!\mathit{olarized}}
\def\Complexesstack{\mathcal{C}\!\mathit{omplexes}}
% \Pic is the operator that assigns to X its picard group, usage \Pic(X)
% \Picardstack_{X/B} denotes the Picard stack of X over B
% \Picardfunctor_{X/B} denotes the Picard functor of X over B
\def\Pic{\mathop{\mathrm{Pic}}\nolimits}
\def\Picardstack{\mathcal{P}\!\mathit{ic}}
\def\Picardfunctor{\mathrm{Pic}}
\def\Deformationcategory{\mathcal{D}\!\mathit{ef}}
\usepackage{cleveref}
\crefname{originaldefinition}{original definition}{original definitions}
\Crefname{originaldefinition}{Original definition}{Original definitions}
\makeatletter
\newlabel{sites-definition-point}{{00Y5}{source}{}{sourcepointer-1}{}}
\newlabel{sites-definition-point@cref}{{[originaldefinition][0][]00Y5}{[1][0][]source}{}{}{}}
\makeatother
\crefname{originalequation}{original equation}{original equations}
\Crefname{originalequation}{Original equation}{Original equations}
\makeatletter
\newlabel{sites-equation-stalk}{{04EH}{source}{}{sourcepointer-2}{}}
\newlabel{sites-equation-stalk@cref}{{[originalequation][0][]04EH}{[1][0][]source}{}{}{}}
\makeatother
\crefname{originalproposition}{original proposition}{original propositions}
\Crefname{originalproposition}{Original proposition}{Original propositions}
\makeatletter
\newlabel{sites-proposition-point-limits}{{00YC}{source}{}{sourcepointer-3}{}}
\newlabel{sites-proposition-point-limits@cref}{{[originalproposition][0][]00YC}{[1][0][]source}{}{}{}}
\makeatother
\crefname{originaldefinition}{original definition}{original definitions}
\Crefname{originaldefinition}{Original definition}{Original definitions}
\makeatletter
\newlabel{definition-geometric-point}{{03PO}{source}{}{sourcepointer-4}{}}
\newlabel{definition-geometric-point@cref}{{[originaldefinition][0][]03PO}{[1][0][]source}{}{}{}}
\makeatother
\crefname{originallemma}{original lemma}{original lemmas}
\Crefname{originallemma}{Original lemma}{Original lemmas}
\makeatletter
\newlabel{lemma-geometric-lift-to-cover}{{03PR}{source}{}{sourcepointer-5}{}}
\newlabel{lemma-geometric-lift-to-cover@cref}{{[originallemma][0][]03PR}{[1][0][]source}{}{}{}}
\makeatother
\title{The geometric fibre functor on the small etale site is the stalk at a point of the site}
\author{The Stacks Project Authors}
\date{stacks-project:etale-cohomology.tex:local-snapshot-20261011; GNU FDL 1.2 or later}
\begin{document}
\maketitle
\noindent\textit{Editorial note.} All delivered mathematical text is the author's own, copied from the pinned Stacks Project chapter file etale-cohomology.tex; no statement or proof step was rewritten, reordered or removed.\par\medskip
\noindent\textit{Editorial note.} The preamble is the chapter's own native preamble excerpt used for this source: the Xy-pic setup at lines 14-18 and the theorem declarations with the mathematical macros at lines 154-217 of sources/stacks-etale-cohomology/preamble.tex. The quoted class line and the xr-hyper externaldocument block are not included because the assembler supplies a plain amsart wrapper; no macro was invented or redefined.\par\medskip
\noindent\textit{Editorial note.} The same-chapter reference is located at its \textbackslash{}label line inside the pinned chapter file; the three cross-chapter pointers are located in the project's own tag list staged in the snapshot (sources/stacks-chow/tags.txt), at the verbatim lines carrying their canonical tags.\par\medskip
\noindent\textit{Allowed standard prerequisites (editorial).}\begin{itemize}
\item Point of a site and its associated stalk functor. Sites, Definition sites-definition-point (tag 00Y5) and Sites, Equation sites-equation-stalk (tag 04EH), named in the statement at lines 3383-3385; the criterion for a functor to be a point is Sites, Proposition sites-proposition-point-limits (tag 00YC), invoked at line 3406. Supply the definition and the criterion that has to be verified; the verification is the counted content.
\item Every scheme etale over an algebraically closed field splits as a disjoint union of copies of the field, and lifting of geometric points along etale coverings. Same chapter, Lemma labelled lemma-geometric-lift-to-cover (source line 3305, Stacks tag 03PR), used at proof lines 3391-3392 and again at 3403-3404. Supplies that U over the geometric point is a disjoint union of copies of the point, and that a geometric point of an object lifts to some member of any etale covering; the functoriality and covering computations are retained in the counted span.
\end{itemize}\medskip
\ifdefined\hypertarget\hypertarget{sourcepointer-1}{}\fi
\noindent\textit{Source reference.} Original definition 00Y5: sources/stacks-chow/tags.txt, source lines 1243--1243. Cross-chapter pointer named in the counted statement at line 3383 for the notion of a point of a site; the verification that u defines such a point is retained verbatim in the counted span.\par\medskip
\ifdefined\hypertarget\hypertarget{sourcepointer-2}{}\fi
\noindent\textit{Source reference.} Original equation 04EH: sources/stacks-chow/tags.txt, source lines 5763--5763. Cross-chapter pointer named twice in the counted statement (lines 3385 and 3411) for the displayed formula defining the stalk functor; the comparison of the two colimits is retained in the counted span.\par\medskip
\ifdefined\hypertarget\hypertarget{sourcepointer-3}{}\fi
\noindent\textit{Source reference.} Original proposition 00YC: sources/stacks-chow/tags.txt, source lines 1250--1250. Cross-chapter standard prerequisite invoked once at proof line 3406 as the criterion for a functor to define a point of a site; the three conditions it requires are verified inside the counted span.\par\medskip
\ifdefined\hypertarget\hypertarget{sourcepointer-4}{}\fi
\noindent\textit{Source reference.} Original definition 03PO: sources/stacks-etale-cohomology/pinned-original-source.tex, source lines 3201--3201. Same-chapter pointer invoked at proof line 3398 for the description of a geometric point as a morphism from the spectrum of an algebraically closed field; the identification of the underlying set of the geometric fibre with that collection is retained in the counted span.\par\medskip
\ifdefined\hypertarget\hypertarget{sourcepointer-5}{}\fi
\noindent\textit{Source reference.} Original lemma 03PR: sources/stacks-etale-cohomology/pinned-original-source.tex, source lines 3305--3305. Same-chapter pointer invoked at proof lines 3391-3392 and 3403-3404 for the shape of the geometric fibre and for surjectivity on coverings; both uses are inputs whose consequences are drawn inside the counted span.\par\medskip
\begin{lemma}
\label{lemma-stalk-gives-point}
Let $S$ be a scheme. Let $\overline{s}$ be a geometric point of $S$.
Consider the functor
\begin{align*}
u : S_\etale & \longrightarrow \textit{Sets}, \\
U & \longmapsto
|U_{\overline{s}}|
=
\{\overline{u} \text{ such that }(U, \overline{u})
\text{ is an \'etale neighbourhood of }\overline{s}\}.
\end{align*}
Here $|U_{\overline{s}}|$ denotes the underlying set of the geometric fibre.
Then $u$ defines a point $p$ of the site $S_\etale$
(Sites, Definition \ref{sites-definition-point})
and its associated stalk functor $\mathcal{F} \mapsto \mathcal{F}_p$
(Sites, Equation \ref{sites-equation-stalk})
is the functor $\mathcal{F} \mapsto \mathcal{F}_{\overline{s}}$
defined above.
\end{lemma}

\begin{proof}
In the proof of
Lemma \ref{lemma-geometric-lift-to-cover}
we have seen that the scheme $U_{\overline{s}}$ is a disjoint union of
schemes isomorphic to $\overline{s}$. Thus we can also think of
$|U_{\overline{s}}|$ as the set of geometric points of $U$ lying over
$\overline{s}$, i.e., as the collection of morphisms
$\overline{u} : \overline{s} \to U$ fitting into the diagram of
Definition \ref{definition-geometric-point}.
From this it follows that $u(S)$ is a singleton, and that
$u(U \times_V W) = u(U) \times_{u(V)} u(W)$
whenever $U \to V$ and $W \to V$ are morphisms in $S_\etale$.
And, given a covering $\{U_i \to U\}_{i \in I}$ in $S_\etale$
we see that $\coprod u(U_i) \to u(U)$ is surjective by
Lemma \ref{lemma-geometric-lift-to-cover}.
Hence
Sites, Proposition \ref{sites-proposition-point-limits}
applies, so $p$ is a point of the site $S_\etale$.
Finally, our functor $\mathcal{F} \mapsto \mathcal{F}_{\overline{s}}$
is given by exactly the same colimit as the functor
$\mathcal{F} \mapsto \mathcal{F}_p$ associated to $p$ in
Sites, Equation \ref{sites-equation-stalk}
which proves the final assertion.
\end{proof}

\end{document}
